	\section{Vector Bundles and Riemannian Geometry}
\label{sec:background-geometry}
Due to differential geometry and the theory of vector bundles being areas of mathematics used by pure mathematicians interested in analytical aspects like curvature, by topologists interested in invariants like Euler characteristics, physicists using the spaces as playgrounds for general relativity theories, and many others, there are numerous competing notational conventions. This thesis lies at the tripoint of differential geometry, dynamical systems, and algebraic topology. Hence, there is no a priori preferred notational convention.
In this chapter, we establish the necessary background in the theory of vector bundles and Riemannian manifolds whilst not proving all the basics; we establish the notation used throughout this thesis. We establish manifolds and vector bundles in parallel, as they are interdependent on one another.
\subsection{Differentiable Structures on Topological Manifolds}
We begin by defining a topological manifold. Afterwards, we equip that space with with increasingly more structure and hence reduce the generality until we arrive at the playground of the thesis: 

\begin{definition}[Topological Manifold] \label{def:vb-topologicalManifold}
	A second countable Hausdorff space $M$ is called a \textbf{topological manifold} of dimension $m\in \N$, if it is locally homeomorphic to $\R^m$. To be precise, for all $p\in M$ there exists an open neighbourhood $U\subseteq M$ of $p$, an open set $V\subseteq \R^m$  and a map $\phi:W\to V$ that is a homeomorphism. We call the map $\phi:U\to V$ a \textbf{chart around $p$ on $M$}, and $\phi^{-1}$ a \textbf{local coordinate system around $p$ on $M$}.
\end{definition}

\begin{definition}[Differentiable Manifold] \label{def:vb-differentiableManifold}
	Let $M$ be a topological manifold of dimension $m$. 
	\begin{enumerate}
		\item A \textbf{differentiable atlas of class} $r\in \N\cup\{\infty\}$ is a family of charts $\mathfrak{A}=\left( \phi_i:U_i \to V_i\right)_{i\in I}$ such that
		\begin{enumerate}
			\item $\bigcup_{i\in I}U_i=M$, meaning that $\left( U_i\right)$ is an open covering of $M$.
			\item For every pair $(i,j)\in I^2$, the \textbf{transition function}:
			\begin{align*}
				\phi_{ij}:\phi_j(U_i \cap U_j)&\to \phi_i(U_i \cap U_j)\\
				x                             &\mapsto \left(\phi_i \circ \phi_j^{-1}\right)(x)
			\end{align*}
			is differentiable of class $r$.
		\end{enumerate}
		We call such an atlas a $C^r$-atlas.
		\item Two $C^r$-atlases $\mathfrak{A}$ and $\mathfrak{B}$ are called \textbf{equivalent} if the family $\mathfrak{A}+\mathfrak{B}=\left(\phi_i,\phi_j\right)_{ij}$ is a $C^r$-atlas.
	\end{enumerate}
	A \textbf{differentiable structure of class $r$} on $M$ is an equivalence class $c$ of $C^r$-atlases. 
	For $r=\infty$, we call the pair $(M,c)$ a smooth manifold.
\end{definition}


\begin{corollary}
	Every transition function $\phi_{ij}$, $i,j\in I^2$ of a differentiable atlas $\mathfrak{A}=(\phi_i)_{i\in I}$ is not merely a homeomorphism but also a diffeomorphism. This follows from the fact that $\phi_{ij}^{-1}=\phi_{ji}$ which provides smooth inverses.
\end{corollary}

\begin{definition}[Smooth Functions] \label{def:vb-smoothFunctions}
	Let $(M,c)$ be a differentiable manifold of class $r$ and $U\subseteq M$ open. We call a continuous function 
	\begin{align*}
		f:U \to \R
	\end{align*} \textbf{differentiable of class $r$}, if for any chart (and hence for all) $(\phi_i)_{i\in I}=\mathfrak{A}\in c$ the compositions $f\circ \phi_i^{-1}$ are differentiable of class $r$. For $r=\infty$ we define:
	\begin{align*}
		\mathcal{E}(U)=\{f :  U \to \R \text{ continouos}~|~ f \text{ is differentiable of class }\infty \} . 
	\end{align*}
\end{definition} 


\begin{corollary}
	Let $(M,c)$ be a smooth manifold of dimension $m$ and $U\subseteq M$ be an open subset. With pointwise defined operations, the set $(\mathcal{E}(U),+,\cdot,\circ)$ becomes an $\R-$algebra. Furthermore, $\mathcal{E}$ becomes a sheaf of $\R$-algebras.
\end{corollary}
\begin{proof}
	There is not really a need for a proof. However, it might help to work through the definition of a sheaf. First, $\mathcal{E}$ is a presheaf, where the restriction in the domain of a function gives the needed restriction homomorphism:
	\begin{align*}
		\text{res}^U_V: \mathcal{E}(U) &\to \mathcal{E}(V) \\
		f &\mapsto \restr{f}{V} \, .
	\end{align*} The required properties of a presheaf are trivial. Furthermore, this gives a sheaf as the requirement of locality is trivial for functions and the property of gluing is also trivial for functions, since differentiability is a local property.
\end{proof}
\begin{definition}
	If $p\in M$ is fix, $f\in \mathcal{E}(U)$ and $g\in \mathcal{E}(U')$ such that $p\in U\cap U'$ we say that $f$ and $g$ have the same \textbf{germ in $p$}, if there is another open neighborhood $W\subseteq U\cap U'$ of $p$ such that $\restr{f}{W}=\restr{g}{W}$. This defines an equivalence relation $\sim_p$. An equivalence class $s$ of local functions around $p$ is called a \textbf{germ in $p$}. We write $s=f_p$, if $s=[f]$ with $f\in \mathcal{E}(U)$. 
	We write 
	\begin{align*}
		\mathcal{E}_p(M)=\left. \left( \sum_{U \text{open},p\in U }\mathcal{E}(U)\right)\middle/ \sim_p \right. \, 
	\end{align*} for the set of germs, and call it the \textbf{stalk in $p$}. Where $\sum$ denotes the co-product (disjoint union) in $\Top$.
\end{definition}
\begin{corollary}
	For a smooth manifold $(M,c)$ the set $\mathcal{E}_p(M)$ inherits an $\R-$algebra structure from the $\mathcal{E}(U)$. Furthermore, it carries a natural (evaluation-)homomorphism:
	\begin{align*}
		\mathrm{eval}_p: \mathcal{E}_p(M) & \to \R \\
		f_p  & \mapsto f(p)\eqcolon f_p(p)
	\end{align*}
	The stalks are also local rings with maximal ideal $\mathfrak{m}_p=\ker(\mathrm{eval}_p)$. Hence, the pair $(M,\mathcal{E})$ gives us a locally ringed space.
\end{corollary}
\begin{proof}
	We only need to verify the local ring structure of the stalks. An element $f_p\in \mathcal{E}_p(M)$ is invertible if and only if $f(p)\neq 0$, equivalently if $f_p\notin \ker(\mathrm{eval}_p)$. Thus $\mathcal{E}_p(M) \big/ \ker(\mathrm{eval}_p)$ is a field, which proves that $\ker(\mathrm{eval}_p)$ is maximal.
\end{proof}



\begin{definition}
	Let $(M,c)$ be a smooth manifold of dimension $m$ and $p\in M$. We call an $\R$-linear map $\delta:\mathcal{E}_p(M)\to \R$ a \textbf{derivation}, if it satisfies the Leibniz rule:
	\begin{align*}
		\delta(f_p\cdot g_p)=\delta(f_p)g_p(p)+f_p(p)\delta(g_p) \quad \text{for all }f_p,g_p\in \mathcal{E}_p(M) \, .
	\end{align*}
	We call $\mathrm{Der}_{\R}(\mathcal{E}_p(M),\R)$ the set of derivations and give it an $\R$-vector space structure by pointwise operations. We define the \textbf{tangent space of $M$ at $p$} to be the vector space
	\begin{align*}
		T_pM \coloneq \mathrm{Der}_{\R}(\mathcal{E}_p(M),\R) \, .
	\end{align*}
\end{definition}
\begin{corollary}
	Let $(M,c)$ be a smooth manifold of dimension $m$ and $\phi:U\to V$ be a chart around $p$ with $x_0=\phi(p)$ (where $\phi \in \mathfrak{A}\in c$). Then we define the derivation
	\begin{align*}
		\partial_i(p) : \mathcal{E}_p(M) &\to \R \\
		f_p & \mapsto \partial_i(p)(f_p)\coloneq \frac{\partial(f\circ \phi^{-1})}{\partial x_i}\big|_{x_0}\,.
	\end{align*}This is well-defined and hence gives a tangent vector. Moreover, the family
	\begin{align*}
		\left(\partial_1(p),\dots,\partial_m(p)\right)
	\end{align*} defines a basis of $T_pM$. Hence, the dimension of $T_pM$ is $m$.
\end{corollary}

This is a good point to take a break from smooth manifolds and introduce the structure of vector bundles. 






%----------Vector bundles-------------------

Analogously to the definition of a smiith manifold, we start by defining families of vector spaces and equip them with increasingly more structure:


\begin{definition}[Vector Bundles]\label{def:vb-vectorBundles}
	A \textbf{complex vector bundle} of rank $m$ is a surjection 
	\begin{equation*}
		\pi: E\to X
	\end{equation*}such that $E$ and $X$ are topological spaces and $\pi$ is continuous. Furthermore $\pi$ satisfies:
	 \begin{enumerate}
	 	\item For all $x\in X$, $\pi^{-1}(x)$ carries the structure of an $m$-dimensional $\C$-vector space.
	 	\item For all $x\in X$, there exist a neighbourhood $U\subseteq X$ and a homeomorphism 
		 	\begin{equation*}
		 		\psi: \restr{E}{U}\coloneq \pi^{-1}(U) \to U\times \C^m \,,
		 	\end{equation*}
		 	that is compatible in the following sense: Let $\pi_1$ denote the projection $ U\times \C^m \to U$ onto the first factor. Then 
		 	\begin{equation*}
		 		\pi_1 \circ \psi=\pi  
		 	\end{equation*} and 
		 	\begin{equation*}
		 		\psi_x: E_x\coloneq \pi^{-1}(x)\to \{x\}\times \C^m 
		 	\end{equation*} is a vector space isomorphism. 
	 \end{enumerate}
	 We call $E$ the \textbf{total space}, $X$ the \textbf{base space}, $E_x$ the \textbf{fibre over $x$} and $\psi$ a \textbf{local trivialization}. 
\end{definition}
\begin{definition}
	Let $\pi: E \to X$ be a vector bundle. A continuous function $s: X\to E$ such that $\pi\circ s=\id_X$ is called a (global) \textbf{section} (if $s:U\subseteq X \to E$ we call it a local section). We denote the space of global sections by $\Gamma(E)$ and by $\Gamma(U)$ the space of local sections over an open subset $U\subseteq X$. 
\end{definition}

\begin{definition}[Homomorphism of Vector Bundles]\label{def:vb-homomorphismOfVectorBundles}
	A \textbf{homomorphism between two vector bundles} $\pi: E\to X$ and $\tilde{\pi}: F\to X$ over the same base space is a continuous function $\phi: E \to F$ such that: 
	\begin{enumerate}
		\item \begin{tikzcd}
			E \arrow[d,"\pi"] \arrow[r,"\phi"]& F \arrow[ld,"\tilde{\pi}"]\\
			X&&
		\end{tikzcd} commutes.
		\item For all $x\in X$ the function $\phi_x: E_x\to F_x$ is linear. 
	\end{enumerate} 
\end{definition}

\begin{definition}[Pullback of Vector Bundles]
	Let $\pi: E\to Y$ be a vector bundle, $X$ be a topological space and $f: X\to Y$ be continuous. Then we define the \textbf{pullback} $f^*{\pi}:f^*(E) \to X$ as follows: 
	The total space $f^*(E)$ is defined as the subspace of $(x,e)\in X\times E$ such that $f(x)=\pi(e)$, and the map $f^*{\pi}$ is the projection to the second factor restricted to said subspace.
\end{definition}
\begin{example}
	Let $\pi:E\to X$ be a vector bundle. For any $Y\subseteq X$ we have the vector bundle 
	\begin{equation*}
		\restr{\pi}{Y}:\pi^{-1}(Y)\to Y\,,
	\end{equation*} which is defined as the pullback along the inclusion of $Y$ into $X$.
\end{example}
\begin{comment}
	
	\begin{definition}[Family of Vector Spaces] \label{def:vb-familiyOfVectorSpaces}
		Let $X$ be a topological space. A \textbf{family of vector spaces over $X$} is a topological space $E$ equipped with :
		\begin{enumerate}[(i)]
			\item a continuous function $\pi: E \to X$,
			\item a finite dimensional vector space structure on each fiber $E_x:=\pi^{-1}(x)$ that is compatible with the subspace topology on $E_x\subseteq E$.
		\end{enumerate}
		We call the map $\pi$ the projection, $E$ the total space and $X$ the base space. 
	\end{definition}
	
	\begin{definition}[Sections]
		Let $\pi: E\to X$ be a family of vector spaces. A function $f:X\to E$ such that for all $x\in X$ we have $\pi\circ f (x)=x$ is called a \textbf{section of $E$}. We denote the space of all sections by $\Gamma(E)$. 
	\end{definition}
	\begin{remark} Let $\pi: E\to X$ be a family of vector spaces.
		By locality of the definition, for all open $U\subseteq X$ the restriction $\restr{\pi}{U}:\pi^{-1}(U)\to U$ defines a new family of vector spaces. With this we can also have sections of the restriced bundle and the collection f those is denoted by $\Gamma(U)$
	\end{remark}















\begin{definition}[Homomorphism of families of vector spaces]\label{def:vb-homomorphismOfFamiliesOfVectorSpaces}
	A \textbf{homomorphism} form one family $\pi:E\to X$ to another family $\Tilde{\pi}:F \to X$ with the same base space is a continuous map $\phi: E\to F$ such that:
	\begin{enumerate}[(i)]
		\item $\Tilde{\pi}\phi=\pi$, i.e. 
		\begin{tikzcd}
			E \arrow[d, "\pi"] \arrow[r, "\phi"] & F \arrow[ld, "\Tilde{\pi}"] \\
			x                                    &                            
		\end{tikzcd}
		\item for each $x\in X$ $\phi_x:E_x\to F_x$ is linear w.r.t. the vector space structure. 
	\end{enumerate} In the case that $\phi$ is bijective (and thereby a point-wise vector space isomorphism) and has an continuous inverse we call it an \textbf{isomorphism} and call $E$ and $F$ isomorphic. 
\end{definition}
\end{comment}



\begin{remark}
	Notice that this construction of the pullback $f^*{\pi}:f^*(E) \to X$ is the simplest space to make the square commute: 
\begin{center}
		\begin{tikzcd}
		f^*(E) \arrow[r,"\pi_2"] \arrow[d,"f^*\pi"]& E \arrow[d,"\pi"]\\
		X \arrow[r,"f"] &Y
	\end{tikzcd}
\end{center} The notion of simplicity is captured by the universal property, that for any other vector bundle $\tilde{\pi}:\tilde{E}\to X$ and function $\tilde{f}:\tilde{E}\to E$ there exists a unique $h:\tilde{E}\to f^*(E)$ (which is given by $h(y)\coloneq (\tilde{\pi}(y),\tilde{f}(y))$ such that the following diagram commutes:

\begin{center}
	\begin{tikzcd}
		\tilde{E} \arrow[ddr,"\tilde{\pi}",bend right] \arrow[dr,"h"] \arrow[drr,"\tilde{f}",bend left] & & \\
		& 	f^*(E) \arrow[r,"\pi_2"] \arrow[d,"f^*\pi"]& E \arrow[d,"\pi"]\\
		&    X \arrow[r,"f"] &Y
	\end{tikzcd}
\end{center}
\end{remark}
\begin{corollary} \label{cor:vb-pullbackIsFunctoriell}
	Let $f: X\to Y$ and  $g:Y\to Z$ be continuous and let $\pi_i:E_i\to Y$ for $i=1,2$ be vector bundles. Then we have the properties:
	\begin{enumerate}
		\item There is a natural isomorphism $g^*f^*(E_1)\cong (f \circ g )^*(E_1)\,.$
		\item $\id_Y^*(E_1)\cong E_1	 
	\end{enumerate}
\end{corollary}

\begin{remark}
	Let $\pi:E\to X$ be a vector bundle. Then the space of sections $\Gamma(E)$ is a (in general, infinit-dimensional) $\R$-vector space. Furthermore, and more relevant to the purposes of this thesis, we view $\Gamma(E)$ (and $\Gamma(U)$ for open subsets) as a $C(X)$ (or $C(U)$) module.  The latter denotes the ring of continuous functions from $E$ (or $U$) to $\R$. 
\end{remark}

\begin{example}
	For any smooth manifold $M$ of dimension $m$ , we can equip the union 
	\begin{equation*}
		TM\coloneq \bigsqcup_{p\in M} T_pM
	\end{equation*} with a topology, such that $\pi: TM\to M$ defines a vector bundle which we call \textbf{tangent bundle}. For this we first define the trivialisations and then equip $TM$ with a topology to make them homeomorphisms: Let $\phi:U\to V$ be a chart on $M$. If we denote the $i$-th coordinate function of $\phi$ by $\phi_i$ we can define the bijection \begin{equation*}
	\tilde{\phi}:\pi^{-1}(U)\to V\times \R^m,\quad \xi_p\mapsto \left(\phi(p),{\footnotesize
	\begin{pmatrix}
		\xi_p(\phi_1)\\
		\vdots\\
		\xi_p(\phi_m)
	\end{pmatrix}
	}\right)\,.
	\end{equation*} Now we declare a subset $W\subseteq TM$ to be open, if and only if for all charts $\phi$ the set $\tilde{\phi}(W\cap \pi^{-1}(U)) \subset V\times \R^m$ is open.
	We can now extend and give meaning to the notation of the local basis vectors $\partial_i(p)\in T_pM$: For any chart $\phi:U\to V$ we have the local sections 
	\begin{equation*}
		\partial_i: U\to TU,\quad p\mapsto \partial_i(p)\,.
	\end{equation*}
	Furthermore, $TM$ is a $2m$-dimensional manifold. Using a smooth atlas on $M$ yields a smooth atlas on $TM$ making it a smooth manifold.
	Notice that the family $\left(\partial_i\right)_i$ is a basis of the $C(U)$-module $\Gamma(U)$. This follows by expressing any section in local coordinates via $\tilde{\phi}$.
\end{example}
Notice that $TM$ as a bundle satisfies more then just the bundle definition: Taking two different charts induces two different local trivialisations. Transitioning between them is not just fibrewise linear, but also smooth. This is captured in the definition:

Henceforth we suppress the differentiable structure from the notation.



\begin{definition}[The Derivative]
	Let $M$ and $M'$ be smooth manifolds. We call a continuous function $f:M \to M'$ \textbf{smooth}, if for all $\phi\in \mathfrak{A}\in c$ and $\phi'\in \mathfrak{A}'\in c'$ the maps
	\begin{align*}
		\phi'\circ f\circ  \phi^{-1}:V\to V'
	\end{align*} are smooth. A given smooth function induces a smooth map between the tangent bundles as manifolds as follows:
	\begin{align*}
		\diff f: TM &\to TM' ~,\diff f_p(\xi_p)(g_p)=\xi((g\circ f)_p)
	\end{align*}Here, $\xi_p\in T_pM$, $g_p\in \mathcal{E}_p(M')$. The function $\diff f$ is called the \textbf{derivative} or \textbf{pushforward} of $f$.
\end{definition}
\begin{definition}[Smooth Vector Bundles]
	A \textbf{smooth vector bundle of rank $r$ } is a surjection $\pi: E \to M$ where the codomain (or base space) is a smooth manifold together with a family of continuous functions (or local trivialisations )
\begin{equation*}
	\Phi_i: \pi{-1}(U_i)\to U_i\times \R^r
\end{equation*} such that  $\bigcup_i U_i=M$ and for all overlapping $U_i\cap U_j\neq \emptyset$, the transition functions
\begin{equation*}
	\Phi_j\circ \Phi_i^{-1}: \pi^{-1}(U_i\cap U_j) \to (U_i\cap U_j)\times \R^r
\end{equation*} are of the form 
\begin{equation*}
	\Phi_j\circ \Phi_i^{-1}(p,v)=(p,g_{i,j}(v)),\quad \text{ where }\quad g_{i,j}: (U_i\cap U_j) \to \GL(r,\R)
\end{equation*} is smooth.
\end{definition}
\begin{corollary}
A smooth vector bundle is a smooth manifold and also a vector bundle. The latter is ensured by the codomain of $g_{i,j}$ which is also called the \textbf{structure group}. We can enforce structure on the bundle by reducing the structure group. As a preview: reducing it to $\GL^+(k,\R)$ the linear maps with positive determinant means that we can induce the concept of orientations. Reducing it to $\O(k)$ lets us equip it with a Riemannian metric. 
\end{corollary}

\begin{definition}\label{def:vb-framesFromTrivialisations}
Let $\pi:E\to M$ be a smooth vector bundle of rank $r$ and let 
\begin{equation*}
	\Phi: \pi^{-1}(U) \to U\times \R^r
\end{equation*} be a local trivialisation.
Then a section is called smooth if it is a smooth map of manifolds. A family of smooth local sections $(s_1,\dots, s_r)$ with $s_i\in \Gamma(U)$ is called a \textbf{local frame} if they form a basis of the $C(U)$-module $\Gamma(U)$. Such a family is obtained from $\Phi$ by defining $s_i(p)\coloneq \Phi^{-1}(p,e_i)$. Similar, from a local frame we can obtain a local trivialisations by defining:
\begin{equation*}
	\Phi: \pi^{-1}(U) \to U\times \R^r\,,\quad v=\sum_{i=1}^{r} \lambda_i s_i(\pi(v)) \mapsto (\pi(v),\lambda_1,\dots, \lambda_r)\,.
\end{equation*}
\end{definition}
\begin{example}
The local trivialisations $\partial_i$ of the tangent bundle of a smooth manifold $M$ form a local frame and make the tangend bundle into a smooth vector bundle.
\end{example}

\begin{definition}[Compact-open topology]
	Let $X,Y$ be topological spaces. Then the set of continuous maps $C(X,Y)$ can be equipped with the topology that has 
	\begin{equation*}
		\{V(K,U)|K\subseteq X \text{ compact }, U\subseteq Y \text{ open}\} \text{ where }V(K,U):=\{f\in C(X,Y) |~f(K)\subseteq U\}
	\end{equation*} as a subbasis. This coincides with the topology of $\Hom(V,W)$ as a subset of $\C^n$ after identification with the corresponding matrix.
\end{definition}


\begin{definition}[Homomorphism Bundle]\label{def:vb-homomorphismBundle}
	Let $W,V$ be vector spaces and $E=V\times X$, $F=W\times X$ the product bundles over $X$. Now any bundle homomorphism $\phi:E\to F$ determines a map $\Phi:X\to \hom(V, W)$ by sending $x\mapsto \phi_x$. With respect to the compact-open topology, this $\Phi$ is continuous. Furthermore, a continuous map $\Phi:X \to \hom(V,W)$ determines a bundle homomorphism $\phi:E\to F$.
\end{definition}

\begin{proof}
	Let $\{e_i\}$ and $\{f_i\}$ be bases of $V$ and $W$. Then each $\Phi(x)$ is a matrix $\Phi_{i,j}(x)$, such that:
	\begin{equation*}
		\Phi(x)e_i=\sum_j\Phi(x)_{i,j}f_j\,.
	\end{equation*} Now $\Phi$ is continuous if and only if $\Phi_{i,j}$ is continuous. Since $\phi(x,v)=(x,\Phi(x)v)$ this is equivalent to $\phi$ being continuous. 
\end{proof}
\begin{lemma}[Isomorphic Vector Bundles]\label{thm:vb-isomorphicVectorBundles}
	Let $E$ and $F$ be vector bundles and $\phi:E\to F$ be a homomorphism. Then $\phi$ is an isomorphism if and only if all $\phi_x$ are isomorphisms. 
\end{lemma}
\begin{proof}
	To show that it is an isomorphism we need an continuous inverse, equivalently, we need to show that $\phi$ is a homeomorphism. First of all we can deduce, that $\phi$ is bijective:
	Assume it wasn't injective, then there would be a fibre in which $\phi_x$ is not injective. Assume that it wasn't surjective. Then there would be a $w\in F$ that has no preimage which means that in the fiber containing $w$ there is no isomorphism. Thus, there is a unique inverse and it remains to verify that it is continuous, which is a local condition. Therefore we can assume that $E$ and $F$ are product bundles. Then we have a continuous map $\Phi:X\to \Hom(V,W)$. Furthermore, since all $\phi_x$ are isomorphisms, we can restrict the image to $\Iso(V,W)$ which is an open subset of $\Hom(V,W)$. Now we can define the map $x\mapsto \Phi(x)^{-1}$ which is continuous since taking inverses is continuous. By the corollary above, we obtain a continuous map $\psi:F\to E$ by sending $(x,w)\mapsto (x,\Phi(x)^{-1}(w))$, which is the inverse. 
\end{proof}
\begin{example}
	Let $M$ be a smooth manifold of dimension $m$, $\phi:U\to V$ a chart and $\tilde{\phi}:TU\to V\times \R^m$ be the corresponding bundle chart. Then it follows from the lemma, that $TU$ is isomorphic to $V\times \R^m$. This is verified by observing that $\partial_i(p)$ is a basis of $T_pU$ for all $p\in U$. In this case $\tilde{\phi}\left(\partial_i\right)=e_i$ giving us isomorphisms in each fiber. 
\end{example}
Now we want to construct new bundles from old ones. There are many operations to construct new vector spaces such as the tensor product, direct sum, $\Hom$-Spaces and many more. Instead of individually proving that those operations applied fiberwise yield new bundles, we subsume these operations within the following definition: 

\begin{definition}[Continuous Functor]
	Let $T$ be a covariant endofunctor in the category of finite-dimensional vector spaces. We call $T$ a \textbf{continuous functor}, if for all $V,W$ the map $T:\hom(V,W)\to \hom(T(V),T(W))$ is continuous. All vector spaces are given the topology induced by their isomorphism to real Euclidean space. If $E$  is a vector bundle we define the set:
	\begin{equation*}
		T(E):=\bigcup_{x\in X}T(E_x)\,.
	\end{equation*} 
\end{definition}
Now we want to equip this with a topology such that $T(E)$ is a bundle over the same base space as $E$.
\begin{definition}[Constructing new bundles]\label{def:vb-constructingNewBundles}
	We start by assuming that $E = X\times V$. In that case we define the topology on $T(E)$ as the product topology $X\times T(V)$. Now suppose that $F = X \times W$ is another product bundle and $f:E\to F$ is a homomorphism. Let $\Phi: X \to \hom(V,W)$ be the corresponding map. Then $T\Phi: X\to \hom(T(V),T(W))$ is continuous by assumption. Thereby, $T(\phi):X\times T(V)\to X\times T(W)$ is continuous. If $\phi$ is an isomorphism, then $T(\phi)$ is also an isomorphism. That follows because it is continuous and an isomorphism fibrewise.
	
	Now assume that $E$ is trivial. Then we can choose an isomorphism $\alpha: E\to X\times V$ and induce a topology on $T(E)$ by pulling it back via $T(\alpha) : T(E) \to X\times T(V)$. This topology does not depend on $\alpha$ by the above. 
	For general vector bundles we do this construction locally. 
	
	If $f:X\to Y$ is continuous, there is a natural isomorphism
	\begin{equation*}
		f^*T(E)\cong Tf^*(E)
	\end{equation*}

\end{definition}
	This construction is carried out in detail by Atiyah in \cite[Chapter 1.2]{atiyah1989k} and can be generalised for continuous functors with several variables both co- and contravariant. Hence this procedure gives us for vector bundles $E$ and $F$ over the same base space the new bundles:
	\begin{enumerate}[(i)]
		\item $E\oplus F $, their direct sum
		\item $E\otimes F$, their tensor product
		\item $\Hom(E,F)$
		\item $E^*$, the dual bundle of $E$ 
		\item $\Lambda^i(E)$ the $i$-th exterior power
	\end{enumerate} and by the natural isomorphism we have natural isomorphisms: 
	\begin{enumerate}[(i)]
		\item $E\oplus F \cong F\oplus E$
		\item $E\otimes F\cong F \otimes E$
		\item $E\otimes (F\oplus F')\cong E\otimes F \oplus E\otimes F'$
		\item $\Hom(E,F) \cong E^*\otimes F$
	\end{enumerate}
	
	
	\begin{corollary}\label{cor:vb-sectionsIntoHom}
		Let $\pi_E:E\to X$ and $\pi_F:F\to X$ be two vector bundles over $X$.
		 Then, the space $\Gamma \left(\Hom(E,F)\right)$ can be identified with the space of bundle morphisms $\hom(E,F)$ since for any section $s:X\to \Hom(E,F)$ we can define the bundle homomorphism $\tilde{s}:E\to F$ by sending $v\mapsto s_{\pi_E(v)}(v)$. 
		 This is a homomorphism since $\pi_F\left( \tilde{s}(v)\right)=\pi_E(v)$\,. 
	\end{corollary}
	\begin{example}[Cotangent bundle]
		Let $M$ be a smooth manifold of dimension $m$. Then we can define the \textbf{cotangent bundle} $T^*M$ to be the dual of the tangent bundle. A local frame of said bundle is given by the fiberwise dual of the frame $(\partial_i) $. Sections into the tangent bundle are called \textbf{vector fields} and sections into the cotangent bundle are called \textbf{one-forms}.  
	\end{example} 
	
	
	\begin{definition}\label{def:vb-linearIdentification}
	Let $V$ be a linear manifold (i.e. a finite-dimensional vector space). Then for any $v\in V$ we have the canonical maps
	\begin{align*}
		R_v:V&\to T_vV\,\\
		x&\mapsto \left(f_v\mapsto \restr{\frac{\partial }{\partial t}}{t=0}f(v+tx)\right)
	\end{align*}
	Let $L:V\to V$ be a linear map viewed as a smooth map between manifolds. Then we have the identity:
	\begin{equation*}
		{\diff L}_{v} \circ R_v =R_{L(v)}\circ L\,.
	\end{equation*} This is proved by a short calculation:
	\begin{equation*}
		\Big({\diff L}_{v} \big( R_v(x)\big)\Big)(f_{L(v)})=R_{v}(x)\left(f\circ L\right)=\restr{\frac{\partial }{\partial t}}{t=0}f\circ L(v+tx)=\underbrace{\restr{\frac{\partial }{\partial t}}{t=0}f\big( L(v)+tL(x)\big)}_{=R_{L(v)}\circ L(x)}
	\end{equation*} Now let $v,v'\in V$. Then we have the translation
	\begin{equation*}
		tr_{v,v'}: V\to V;\quad x\mapsto (x+(v'-v))
	\end{equation*} This allows us to relate $R_v$ and $R_{v'}$ as follows:
	\begin{equation*}
		\left(\restr{\diff tr_{v,v'}}{v}(R_v(x))\right)(f_{v'})=\restr{\frac{\partial }{\partial t}}{t=0}\left(f\circ tr_{v,v'}(v+tx)\right)=\restr{\frac{\partial }{\partial t}}{t=0}\left(f(v'+tx)\right)=R_{v'}(x)(f_{v'}) \,.
	\end{equation*} Hence:
	\begin{equation*}
		R_{v'}=\diff tr_{v,v'}\circ R_v\,.
	\end{equation*} 
	The \textbf{linear identification } is the map $Q$ given on the whole tangent bundle:
	\begin{equation*}
		Q:TV\to V;\quad \xi_p \mapsto R_p^{-1}(\xi_p)\,,
	\end{equation*} and satisfies for any linear $L : v\to V$ 
	\begin{equation*}
		Q\circ \diff L=L\circ Q\,.
	\end{equation*}
\end{definition}

	\begin{corollary}\label{cor:vb-vectorFieldsTakeFunctions}
		Let $f:M\to \R$ be smooth. Then we can interpret $\diff f$ as a one-form. To be precise, let $Q : T\R\to \R$ denote the linear identification with respect to the identity on $\R$. For $v\in \Gamma( TM)$ we have
		\begin{align*}
			(Q\circ \diff f)(v)=v(f)
		\end{align*} Here on the right side, $f$ denotes a map $p\mapsto f_p$ such that $v(f)(p)\coloneq v_p(f_p)$ is well defined. We retain this notation so that vector fields can take functions as an input.
	\end{corollary}

		\begin{proof}Let $\phi:U \to V$ be a chart of $M$. 
			We prove this by showing it for the local sections $\partial_i$ and the general case follows by $C(U)$-linearity, since those sections form a local frame. Now let $g_p\in \mathcal{E}_p(\R)$ and $\phi(p)=x_0$. Then
		\begin{equation*}
		\diff f _p(\partial_i(p))(g_p)= 
		\partial_i((g\circ f)_p)=
		=\restr{\frac{\partial }{\partial x_i}}{x_0}(g\circ f\circ \phi^{-1})
		=\frac{\partial}{\partial_x}_{f_p(p)}(g_p) \cdot  \partial_i(p)(f_p)
		\end{equation*}
		Hence, applying $Q$ to the calculation yields $(Q\circ \diff f)(v)=v(f)$\,.
	\end{proof}
	
	
\begin{definition}[A Metric]
	Let $M$ be a smooth manifold.
	A section $g\in \Gamma(T^*M\otimes T^*M)$ into the tensor product of the dual tangent bundle with itself is a \textbf{Riemannian metric} if the following are satisfied for all $p\in M$:
	\begin{enumerate}
		\item $g$ is \textbf{non-degenerate}, meaning if for a $v_p\in T_pM$ we have $g_p(v_p,w_p)=0$ for all $w_p\in T_pM$ then $v_p=0$.
		\item $g$ is \textbf{symmetric}, meaning $g_p(v_p,w_p)=g_p(w_p,v_p)$ for all $ v_p,w_p\in T_pM$ .
		\item $g$ is \textbf{positive definite}, meaning that $g_p(v_p,v_p)\geq 0 $ for all $v_p$ with equality only for $v_p=0$.
	\end{enumerate} We call a tuple $(M,g)$ a \textbf{Riemannian manifold}.
\end{definition}
\begin{remark}
	Equipping a smooth manifold $M$ of dimension $m$ with a Riemannian metric is equivalent to reducing the structure group of the tangent bundle from $\GL(m)$ to $\O(m)\,.$
	To be precise: A Riemannian metric $g$ lets us canonically choose an atlas with structure group $\O(m)$ and an atlas with such a structure group lets us canonically choose a Riemannian metric. 
\end{remark}
\begin{proof}
	We start by assuming that we have an atlas with structure group $\O(m)$. In Definition \ref{def:vb-framesFromTrivialisations} we saw that a covering of local trivializations is equivalent to a covering of smooth frames. Now in a single trivialization 
	\begin{equation*}
		\Phi: \pi^{-1}(U)\eqcolon TU \to U\times \R^m
	\end{equation*} we call the projection to the $\R^m$-part (the second factor) $\pi_2$ and define 
\begin{align*}
	g: U &\to T^*U \otimes T^*U \\
	p&\mapsto  \begin{array}{rl}
		g_p^\Phi: T_pU\times T_pU &\to \R \\
		(\nu_p,\xi_p)    &\mapsto g_p^\Phi(\nu_p,\xi_p)\coloneq \pi_2\left(\Phi(\nu_p)\right)~\cdot~ \pi_2 \left(\Phi(\xi_p)\right)\,,
	\end{array} 
\end{align*} where $\cdot$ denotes the standart scalar product. This gives us locally a Riemannian metric. If we do this for all trivializations we have to make sure, that the result is well defined. So assume that two trivializations $\Phi$ and $\Psi$ overlap in $V$. The structure group is $\O(m)$, $\Phi=R\Psi$ for some function 
\begin{equation*}
	R: V\to \O(m)\,.
\end{equation*} Let $p$ be in said overlapping and $(\nu_p,\xi_p)$ be from $T_pM\times T_pM$. Then since $R$ is orthogonal we can calculate the independence from the trivialization:
\begin{equation*}
	g_p^\Phi(\nu_p,\xi_p)=\pi_2\left(R\Psi(\nu_p)\right)\,\cdot \, \pi_2 \left(R\Psi(\xi_p)\right)=\pi_2\left(\Psi(\nu_p)\right)~\cdot~ \pi_2 \left(\Psi(\xi_p)\right)=	g_p^\Psi(\nu_p,\xi_p)\,.
\end{equation*}

To get an orthonormal atlas from a smooth one given a Riemannian metric, we start by building a covering of local frames from the trivialisations. We then apply Gram-Schmidt to the local frames to get orthonormal frames. Now a covering of orthonormal frames gives a covering of trivializations, where the structure group consists of basis changes between orthonormal bases and hence is $\O(m)$.
\end{proof}

\begin{definition}[The Gradient]
	Assume that $(M,g)$ is a smooth manifold together with a Riemannian metric. Let $f:M\to \R$ be a smooth function and $Q: T\R\to \R$ the linear identification. We define its \textbf{gradient} to be a vector field $\grad(f)\in \Gamma(TM)$ such that for any vector field $V$:
	\begin{align*}
		 Q\circ \diff f (V) = g(\grad(f),V) \, .
	\end{align*} Notice how the non-degeneracy of $g$ ensures existence and uniqueness. 
\end{definition}
\begin{definition}\label{def:vb-localDescriptionOfGradientGij}
	Let $(M,g)$ be a Riemannian manifold and $\phi:U\to V$ be a chart. We define the smooth functions $g_{ij}:U\to \R$ to be:
	\begin{align*}
		g_{ij}(p)=g_p(\partial_i(p),\partial_j(p)) \, .
	\end{align*}
	Then if two vector fields $v,w$ over $U$ are given with $v=\sum v_i \partial_i$ and $w=\sum _i w_i \partial_i$ we can calculate the metric locally:
	\begin{align*}
		g(v,w)=\sum_{i,j}g_{ij} v_iw_j\, :U\to \R
	\end{align*} Furthermore, we can invert the matrices $(g_{ij}(p))_{ij}$ for all $p$ (which is a smooth procedure as can be seen by the construction via Cramer's rule) to get smooth functions 
	\begin{align*}
		g^{ij}: U&\to \R\\
		p&\mapsto g^{ij}(p)= \left( \big(g_{kl}(p)\big)_{kl}^{-1} \right)_{ij}
	\end{align*}
\end{definition}
\begin{remark}
	Due to the rare appearance of big summations with vectors and co-vectors in this work, we omit the Einstein summation convention and assign all coordinate entries lower indices. 
\end{remark}



\begin{lemma}[Local Description of the Gradient]\label{lem:vb-localDescriptionOfGradient}
	Let $(M,g)$ be a smooth Riemaniannian manifold of dimension $m$ and $\phi:U\to V$ be a chart. Then the gradient has the local form:
	\begin{align*}
		\grad(f) =\sum_{i,j}g^{ij}  \left(\partial_i (f) \right)\cdot  \partial_j \, .
	\end{align*}
\end{lemma}

\begin{proof} Assume that $v,w\in \Gamma( TU)$ such that $v=\sum_i v_i\partial_i$ and $w=\sum_i w_i\partial _i$ (Here $v_i$ and $w_i$ are smooth functions $U\to \R$). Then  $g(v,w)=\sum_{i,j} g_{ij} v_i w_j$.
	Suppose that $\grad(f)=\sum_j G^j \partial_j$ and $Q$ is the linear identification $T\R\to \R$ in each tangent space. Then by definition of the gradient and corollary \ref{cor:vb-vectorFieldsTakeFunctions} we have:
	\begin{align*}
		\partial_j(f)=	Q \circ \diff f (\partial_j)= g(\grad(f),\partial_j)= g(\sum_i G^i \partial_i,\partial_j)=\sum_{i}g_{ij} G^i
	\end{align*} Reading this as a pointwise matrix product:
	\begin{align*}
		&(G^1,\cdots G^m)(g_{ij})=\left(\partial_1 (f),\cdots,\partial_m (f)  \right) \\
		\text{ which gives } &(G^1,\cdots G^m)=\left(\partial_1 (f),\cdots,\partial_m (f)  \right) (g^{ij})\, .
	\end{align*}  But then we can conclude that $G^j=\sum_i g^{ij}\partial_i (f)$.
	
\end{proof}

\begin{definition}[A Dynamical System]
	Let $M$ be a smooth manifold. A \textbf{dynamical system on $M$} (sometimes allso calld a \textbf{flow on $M$}) is given by a smooth map $\phi:\Omega \to M$ such that:
	\begin{enumerate}
		\item $\{0\}\times M \subseteq \Omega \subseteq \R \times M$ open, and for any $p\in M$ the set $I(p)\coloneq \pi_1(\Omega\cap \R\times \{p\})$ is connected. ($\pi_1$ is the projection onto the first factor)
		\item For all $p\in M$ and $t\in I(p)$ we have that $s\in I(\phi(t,p))$ if and only if $s+t\in I(p)$. Then we have:
		\begin{align*}
			\phi(s,(\phi(t,p)))=\phi(s+t,p)
		\end{align*} 
		\item For all $p$ we require $\phi(0,p)=p$. 
	\end{enumerate}
	We will write $\phi_t(p)\coloneq \phi(t,p)$ to keep the notation manageable. 
	If for all $p$ $I(p)=\R$ we call the system \textbf{global}. 
\end{definition}

\begin{corollary}
	For a global dynamical system $\phi$ on a smooth manifold $M$ the map
	\begin{align*}
		\rho: \R \to \Diff(M),\quad t\mapsto \phi_t
	\end{align*}is a homomorphism. To verify well-definedness notice how for all $t\in \R$, $\phi_{-t}$ is an inverse of $\phi_t$. The map $\rho$ is called a \textbf{one-parameter group of diffeomorphisms}.
\end{corollary}
\begin{definition}[Vector Field of a Dynamical System]
	Let $\phi: \Omega \to M$ be a dynamical system on a smooth manifold $M$. For a chart $\psi:U\to V$ of $M$ we have the chart $\id\times \psi$ of $\Omega$. For this chart we denote the associated local frame by $(\partial_t,\partial_1,\dots, \partial_m)$.  Using this we define the mapping
	
	\begin{align*}
		M &\to TM\\
		p& \mapsto X(p)\coloneq \diff \phi (\partial_t(0,p))
	\end{align*} and call it the \textbf{associated vector field}. This is in fact a smooth vector field. 
\end{definition}
\begin{definition}[Dynamical System of a Vector Field]
	Let $X\in \Gamma(TM)$ be a global vector field on the smooth manifold $M$. Then we call a dynamical system $\phi$ on $M$ the \textbf{associated dynamical system} to $X$, if for all $t\in \R$ and $x\in M$ 
	\begin{equation*}
		\diff \phi \left(\partial_t(t,x )\right)=X(\phi_{t}(x))\,.
	\end{equation*} 
\end{definition}
\begin{remark}
	Regarding the left-hand side of the definition $\diff \phi \left(\partial_t(t,x )\right)$ we could also differentiate the map 
	\begin{equation*}
		t \mapsto \phi_t(x)\,,
	\end{equation*} for fixed $x$ at $t$ and evaluate this at the constant vector field $t\mapsto R_t(1)\in T_t \R$ from definition \ref{def:vb-linearIdentification}. 
\end{remark}
These two constructions are mutually inverse, giving the following theorem:
\begin{theorem}
	Let $M$ be a smooth manifold. Then there is a one-to-one correspondence between global dynamical systems and vector fields given by the constructions above.
\end{theorem}

